% Pursuing Stacks (cotes 134-2 to 134-8): archive cote 134-4, site page 22 (PDF page 23), typescript page 110.
% typescript. Draft compared: ps3.tex:1079-1138.
% AI-transcribed by Claude Opus 5.5 (claude-opus-5-5, adaptive thinking, effort high), Message Batches API, batch ps2 (msgbatch_01NVwr1Yuo8qJf5rykeGBXPK), 2026-11-12; not yet checked by eye.
\scanpage{134-4/22}
\mspage{110}
the way it is handled),  and the notion of \uline{contractible objects}. Another
impression was about the dependance of these notions \over{on}{upon} a preliminary no-%
tion of "weak equivalence"  namely upon $W\hand{\subset} F\!\!/(M)$ \over{is}{being} a rather loose one,
Thus the construction of homotopy notions in terms of a given interval\xout{x}
$\underline{I}$ (including the category $M_{\underline{I}}$ with the canonical functor $M\hand{\to}M_{\underline{I}}$ ) is valid
for any interval in any category M with final object and binary products
(instead of binary products,  it is even enough that I be``squarable''\hand{,} name-%
ly all products XxI exist in M).  As for the W-homotopy notions,  they de-%
pend on W via the corresponding notion of W-asphericity over e,  which is
\xout{thx} at first sight \uline{the} natural condition to impose upon an interval $\underline{I}$ ,
in order for the corresponding $\underline{I}$-\unsure{homotopy} notions to fit nicely with W
(as expressed in the homotopy lemma).  But then we noticed that a much
weaker condition than asphericity on $\underline{I}$ suffices  -  namely that the two
sections $\delta_o,\delta_1$ of I over e be W-homotopic,  which means essentially \inserted[typed]{the ``points''} \deleted{they}
\xout{\unsure{$\delta$IxI}}
\xout{\unsure{$\delta$IxI}}
\xout{canxbexjoinedxxby} of I they define can be ``joined'' by a finite chain of
W-aspheric intervals mapping into I.  This strongly suggests (in view of
the main application we have in mind,  namely to the study of modelizers)
that the natural condition to impose upon intervals,  in most contexts of
interest to us,  will be merely 0-connectedness. But this notion is \uline{intrin}
-\uline{sic to the category} M,  irrespective again of \hand{t}he choice of any W;  and
therefore the corresponding homotopy notions in M will turn out \over{i}{(}in the
cases at least of greatest interest to us)  to be equally intrinsic to the
category M.   On the other hand (and here fits in the third main impressi-%
on that peeled out two days ago),  the work carried through so far in
view of the ``observation'' and the (naive) ``question'' of last week (pages
94 and 95) strongly suggests that \inserted{in} the nicest modelizers (including  (Cat)
and the elementary modelizers,  presumably),  the notion of weak equivalen
-ce \inserted{W} can be described in terms of the homotopy notions,  more specifically
in terms of the notion of contractib\hand{le} \deleted{elements} objects  \hand{(} when exactly
and how should appear in due course\hand{)}. Th\hand{u}s it will follow that for those
modelizers,  the modelizing structure W itself is uniquely determined in
terms of the intrinsic category structure  -  thus any equivalence b\hand{e}t-%
ween the underlying categories of any two such modelizers should autom\hand{at}i-%
cally be model-preserving !  It will be rather natural to call  the modeli
-zers which fit into this idyllic picture \uline{canonical modelizers} ,  as their
modelizing structure W is indeed canonically determined by the category
structure.  Next thing then would be to try to gain an overall view of
how to get\xout{x} ``all''  canonical modelizers,  if possible in as concrete terms
as the overall view we got upon elementary modelizers A\^{} in terms of \inserted{the corresponding} test
categories A .

\margin{22}
